\chapter{The ring \texorpdfstring{\(\mathbb Z[\phi]\)}{Z[phi]} and the 120 vertices}
\label{ch:ArchonFull600Theory_Geometry}

This chapter is Milestone 2 of the handoff: O1, O2, Lemma 1 (counts, norms
and the three families, with the two external results E1 and E2 as stated
hypotheses), the supporting-facet test and the completeness of the facet
list (FC), and the combinatorial part of Lemma 3 (N1). Source line numbers
refer to \texttt{research/Full\_600cell\_Puzzle\_Theory.tex}.

\section{The exact ring}

\begin{definition}[Golden ratio identities]
  \label{def:golden-ratio}
  \lean{ArchonFull600Theory.goldenRatioIdentities}
  \leanok
  % SOURCE: Definitions and mathematical sources, lines 273--277
  % SOURCE QUOTE: \phi=\frac{1+\sqrt5}{2},\qquad \tau=\phi^{-1}=\phi-1.
  % \end{equation} Thus \begin{equation}
  % \phi^2=\phi+1,\qquad \tau^2=2-\phi.
  \textit{Source: Definitions (O1).}
  Put \(\phi=(1+\sqrt5)/2\) and \(\tau=\phi^{-1}\). Then
  \(\phi^2=\phi+1\), \(\tau=\phi-1\) and \(\tau^2=2-\phi\).
\end{definition}

\begin{lemma}[Pair arithmetic]
  \label{lem:pair-arithmetic}
  \lean{ArchonFull600Theory.ZPhi.pairArithmetic, ArchonFull600Theory.ZPhi.toReal_mul,
    ArchonFull600Theory.ZPhi.toReal_injective, ArchonFull600Theory.ZPhi.positive_iff}
  \leanok
  % SOURCE: The finite checker and its scope, lines 616--629
  % SOURCE QUOTE: It represents an element by \((a,b)\) and uses
  % \begin{equation}
  % (a+b\phi)(c+d\phi)
  % =(ac+bd)+(ad+bc+bd)\phi.
  % \end{equation} Its sign test compares integers after writing
  % \begin{equation}
  % 2(a+b\phi)=(2a+b)+b\sqrt5,
  % \end{equation} with sign cases handled before squaring.
  \textit{Source: Finite checker (O1).}
  Represent \(a+b\phi\in\mathbb Z[\phi]\) by the integer pair \((a,b)\).
  The product \((a,b)(c,d)=(ac+bd,\ ad+bc+bd)\) is the product of
  \(\mathbb Z[\phi]\); distinct pairs are distinct real numbers; and the
  sign test, which writes \(2(a+b\phi)=(2a+b)+b\sqrt5\) and handles the sign
  cases before squaring, accepts a pair exactly when its real value is
  positive.
\end{lemma}
\begin{proof}
  \leanok
  The product formula is \(\phi^2=\phi+1\). Injectivity is the
  irrationality of \(\phi\). For the sign test, compare \(A\) with
  \(-B\sqrt5\) case by case on the signs of \(A\) and \(B\), squaring only
  when both sides are nonnegative.
\end{proof}

\section{The vertex set}

\begin{definition}[The vertex set \(V\)]
  \label{def:vertex-set}
  \lean{ArchonFull600Theory.vertexSet, ArchonFull600Theory.vertexTable,
    ArchonFull600Theory.range_vertexR}
  \leanok
  % SOURCE: Lemma 1, lines 344--372; the problem's coordinate families, through the
  % even permutation taking (phi,1,tau,0) to (0,1,phi,tau) at lines 362--365
  % SOURCE QUOTE: The permutation taking \begin{equation}
  % (\phi,1,\tau,0)\quad\text{to}\quad(0,1,\phi,\tau)
  % \end{equation} has index sequence \((4,2,1,3)\) and is even.
  \textit{Source: Definitions and Lemma 1 (O2).}
  \(V\subset\mathbb H=\mathbb R^4\) is the union of the three families
  \((\pm2,0,0,0)\), \((\pm1,\pm1,\pm1,\pm1)\) and
  \((0,\pm1,\pm\phi,\pm\tau)\), each with all even permutations of the
  coordinate positions and independent signs. The table
  \texttt{vertexTable} lists the same 120 points with coordinates in
  \(\mathbb Z[\phi]\), in the order of the reference script.
\end{definition}

\begin{lemma}[Coordinate identification]
  \label{lem:coordinate-identification}
  \lean{ArchonFull600Theory.vertexFamilies, ArchonFull600Theory.vertexCount,
    ArchonFull600Theory.vertexNorm, ArchonFull600Theory.halfVertexSet,
    ArchonFull600Theory.CoordinateSources, ArchonFull600Theory.coordinateIdentification}
  \leanok
  \uses{def:vertex-set}
  % SOURCE: Lemma 1, lines 342--383
  % SOURCE QUOTE: The set \(V\) in the problem has \(120\) distinct elements, every
  % element has norm \(2\), and \(V/2\) is a binary icosahedral subgroup of
  % \(\operatorname{Sp}(1)\). Its convex hull is a regular 600-cell.
  \textit{Source: Lemma 1 (L1), with E1 and E2 as hypotheses.}
  The three families have \(8\), \(16\) and \(96\) elements and are
  pairwise disjoint, so \(V\) has \(120\) elements. Every element has norm
  \(2\). \(V/2\) equals the coordinate set of the quaternion coordinate
  result. Hence, given E1 (that coordinate set is the binary icosahedral
  group), E2 (the elements of a binary icosahedral group are the vertices of
  a regular 600-cell) and that multiplication by \(2\) preserves
  regularity, \(V/2\) is binary icosahedral and \(\operatorname{conv}V\) is
  a regular 600-cell.
\end{lemma}
\begin{proof}
  \leanok
  The counts, the disjointness and the norms are finite computations in
  \(\mathbb Z[\phi]\) on the table, transferred to \(\mathbb H\) by the
  injective coordinate map. The identification of \(V/2\) with the
  coordinate set uses the even permutation of index sequence
  \((4,2,1,3)\). The two external results are fields of the structure
  \texttt{CoordinateSources}, whose predicates for ``binary icosahedral''
  and ``regular 600-cell'' are left abstract.
\end{proof}

\section{The edge graph and the face counts}

\begin{definition}[Edge graph]
  \label{def:edge-graph}
  \lean{ArchonFull600Theory.edgeGraph, ArchonFull600Theory.edgeGraph_adj_iff,
    ArchonFull600Theory.inner_vertexR_le}
  \leanok
  \uses{def:vertex-set}
  % SOURCE: The finite checker and its scope, line 639
  % SOURCE QUOTE: Joins pairs with inner product \(2\phi\).
  \textit{Source: Finite checker, step 2.}
  Two vertices are adjacent when their inner product is \(2\phi\). Every
  other pair of distinct vertices has a smaller inner product, by the exact
  sign test on all \(120^2\) pairs.
\end{definition}

\begin{lemma}[Edge length and exposed vertices]
  \label{lem:metric-data}
  \lean{ArchonFull600Theory.exposedVertices, ArchonFull600Theory.edgeLength}
  \leanok
  \uses{def:edge-graph}
  % SOURCE: Lemma 3, lines 448--451
  % SOURCE QUOTE: Every point of \(V\) is an exposed vertex: the functional
  % \(x\mapsto v\cdot x\) has its unique maximum at \(v\) among points of
  % \(V\), since all have norm \(2\). Thus \(P\) has exactly \(120\)
  % vertices and circumradius \(2\).
  \textit{Source: Lemma 3 (L3), metric part.}
  Every point of \(V\) has norm \(2\); the exposed points and the extreme
  points of \(P=\operatorname{conv}V\) are exactly the \(120\) points of
  \(V\). Distinct vertices are at distance at least \(2/\phi\), with
  equality exactly for adjacent vertices.
\end{lemma}
\begin{proof}
  \leanok
  For \(v\in V\), the functional \(x\mapsto v\cdot x\) is at most
  \(2\|x\|\le4\) on \(P\), with equality only at \(x=v\) by the equality
  case of Cauchy--Schwarz. The squared distance of two vertices is
  \(8-2\,v\cdot w\ge8-4\phi=(2/\phi)^2\).
\end{proof}

\begin{theorem}[Face counts and links]
  \label{thm:face-counts}
  \lean{ArchonFull600Theory.faceCounts, ArchonFull600Theory.vertexLinkCounts,
    ArchonFull600Theory.incidenceCounts, ArchonFull600Theory.eulerCharacteristic}
  \leanok
  \uses{def:edge-graph}
  % SOURCE: Lemma 3, lines 437--443 and 487--505; handoff N1
  % SOURCE QUOTE: (f_0,f_1,f_2,f_3)=(120,720,1200,600).
  % \end{equation} Its facets are regular tetrahedra, five facets meet at
  % each edge, twenty meet at each vertex, and its vertex links are
  % icosahedral.
  \textit{Source: Lemma 3 (L3), combinatorial part, and N1.}
  The edge graph has \(120\) vertices, \(720\) edges, \(1200\) triangles
  and \(600\) four-cliques. Every vertex has degree \(12\) and lies in
  \(12\) edges, \(30\) triangles and \(20\) four-cliques, the numbers of
  vertices, edges and triangles of its link. Every edge lies in \(5\)
  triangles and \(5\) four-cliques, and every triangle in exactly \(2\)
  four-cliques. The alternating sum \(120-720+1200-600\) is \(0\).
\end{theorem}
\begin{proof}
  \leanok
  A kernel-checked certificate on the adjacency bitmask gives \(12\) common
  neighbours for every vertex, \(5\) for every edge and \(2\) for every
  triangle. For any finite graph with constant numbers of common
  neighbours \(d\), \(k\), \(m\) of vertices, edges and triangles, double
  counting of nested cliques gives \(2f_1=f_0d\), \(3f_2=f_1k\),
  \(4f_3=f_2m\) and the local counts.
\end{proof}

\section{The supporting-facet test and completeness}

\begin{lemma}[Supporting-facet test]
  \label{lem:supporting-facet}
  \lean{ArchonFull600Theory.supportingFacet, ArchonFull600Theory.facetGram,
    ArchonFull600Theory.finrank_cliqueTetra}
  \leanok
  \uses{thm:face-counts}
  % SOURCE: The finite checker and its scope, lines 654--662
  % SOURCE QUOTE: For a tetrahedron with vertices \(v_1,\ldots,v_4\), the Gram matrix has
  % diagonal entries \(4\) and off-diagonal entries \(2\phi\). It is
  % positive definite. Moreover, \begin{equation}
  % N=\sum_i v_i,\qquad
  % N\cdot v_i=4+6\phi,\qquad
  % \|N\|^2=16+24\phi.
  % \end{equation} Thus the supporting tests certify actual tetrahedral
  % facets, not merely graph cliques.
  \textit{Source: Finite checker, step 4 (FC).}
  Let \(v_1,\ldots,v_4\) be the vertices of a four-clique and
  \(N=\sum_iv_i\). The Gram matrix of the \(v_i\) has diagonal entries
  \(4\) and off-diagonal entries \(2\phi\), and it is positive definite.
  Moreover \(N\cdot v_i=4+6\phi\), \(\|N\|^2=16+24\phi\), and
  \(N\cdot v<4+6\phi\) at every other vertex. So \(N\cdot x\le4+6\phi\) on
  \(P\), with equality exactly on the tetrahedron
  \(\operatorname{conv}\{v_1,\ldots,v_4\}\), whose affine span has
  dimension \(3\).
\end{lemma}
\begin{proof}
  \leanok
  \(N\cdot v_i=4+3\cdot2\phi\) and \(\|N\|^2=4(4+6\phi)\). At another
  vertex each of the four inner products is at most \(2\phi\), and
  \(8\phi<4+6\phi\) since \(\phi<2\). The face of a finite convex hull
  cut out by a bounded functional is the hull of the points where the
  bound is attained. A combination \(\sum c_iv_i\) has squared norm
  \((4-2\phi)\sum c_i^2+2\phi(\sum c_i)^2\), which vanishes only for
  \(c=0\); this gives positive definiteness and affine independence.
\end{proof}

\begin{lemma}[Cliques span exposed faces]
  \label{lem:clique-face}
  \lean{ArchonFull600Theory.cliqueFace}
  \leanok
  \uses{thm:face-counts}
  % SOURCE: Lemma 3, lines 461--467
  % SOURCE QUOTE: Thus the functional \(v+w\) exposes precisely their
  % segment.
  \textit{Source: Lemma 3 (L3), the exposure of edges, extended to all cliques.}
  Let \(s\) be an \(n\)-clique of the edge graph and \(N\) the sum of its
  vertices. Then \(N\cdot x\le4+2(n-1)\phi\) on \(P\), with equality
  exactly on the convex hull of the vertices of \(s\). In particular the
  functional \(v+w\) of an edge exposes precisely its segment, and the
  triangles and four-cliques span exposed faces of \(P\).
\end{lemma}
\begin{proof}
  \leanok
  At a vertex of \(s\) the value is \(4+(n-1)2\phi\); at any other vertex
  it is at most \(n\cdot2\phi\), which is smaller since \(2\phi<4\). The
  face lemma for finite convex hulls finishes the proof.
\end{proof}

\begin{definition}[Facets and the polytope facts E4]
  \label{def:facet-facts}
  \lean{ArchonFull600Theory.IsFacet, ArchonFull600Theory.FacetsAdjacent,
    ArchonFull600Theory.FacetFacts}
  \leanok
  % SOURCE: Lemma 3, lines 494--495; The finite checker and its scope, lines 664--667; handoff E4
  % SOURCE QUOTE: The facet adjacency graph of a convex polytope is connected, so a
  % nonempty such list contains all facets.
  \textit{Source: handoff E4, as a stated hypothesis.}
  A facet of \(P\) is a set \(\{x\in P: l(x)=M\}\) for a linear functional
  \(l\) with \(l\le M\) on \(P\), whose affine span has dimension \(3\).
  Two facets are adjacent when they are distinct and meet in a set of
  affine dimension \(2\). The structure \texttt{FacetFacts} records the two
  standard facts the completeness argument uses: a ridge lies in no facet
  other than the two that meet in it, and the facet adjacency graph is
  connected.
\end{definition}

\begin{theorem}[Completeness of the facet list]
  \label{thm:facet-completeness}
  \lean{ArchonFull600Theory.isFacet_cliqueTetra, ArchonFull600Theory.facetCompleteness}
  \leanok
  \uses{lem:supporting-facet, def:facet-facts}
  % SOURCE: The finite checker and its scope, lines 663--667
  % SOURCE QUOTE: If every triangular ridge of the listed tetrahedra has both its
  % incident facets in the list, the list is closed under facet adjacency.
  % The facet adjacency graph of a convex polytope is connected, so a
  % nonempty such list contains all facets.
  \textit{Source: Finite checker, completeness argument (FC).}
  The tetrahedron of every four-clique is a facet of \(P\). Under the
  polytope facts E4, the facets of \(P\) are exactly these tetrahedra, and
  there are \(600\) of them.
\end{theorem}
\begin{proof}
  \leanok
  Let \(G\) be a facet adjacent to the tetrahedron \(T_s\) of a
  four-clique \(s\). The ridge \(T_s\cap G\) is the hull of the vertices of
  \(s\) where the functional of \(G\) is attained; it has dimension \(2\),
  so it is the triangle of a three-clique \(r\subset s\). That triangle
  lies in exactly two four-cliques, \(s\) and \(s'\), and the ridge lies in
  \(T_{s'}\), so \(G=T_{s'}\) by E4. Since the facet adjacency graph is
  connected and the list is nonempty, every facet is listed. Distinct
  four-cliques give distinct tetrahedra, since a vertex lies in \(T_s\)
  exactly when it belongs to \(s\).
\end{proof}
